\begin{textsample}{NEG Dim 6 – global\_south\_2024\_02 – Score -1.00 – global\_south\_2024\_02/106247796.txt}  \label{ex:f6_neg_277}
What is the significance of the Palestinian Authority government resigning, and what will it effectively mean?
Why are West Bank and Gaza under different administrations?
All you need to know.
Palestinian Prime Minister Mohammad Shtayyeh at a cabinet meeting in Ramallah in the Israeli-occupied West Bank on February 26. ( Photo : Reuters ) With no end to the Gaza war in sight, the West Bank ' s Palestinian Authority government resigned on Monday ( February 27 ).
The move is being seen as the first step towards an overhaul of the Palestinian Authority, so that it can govern Gaza too after the war, an arrangement that the US is pushing.
West Bank, Gaza, and East Jerusalem together make up the state of Palestine.
All these territories are under varying degrees of Israeli occupation.
What is the significance of the West Bank government resigning, and what will it effectively mean?
Can it help shorten the war in Gaza?
We explain.
Advertisement Palestinian Authority, Fatah, and Hamas The Palestinian Authority ( PA ) @ @ @ @ @ @ @ @ @ @ which Israel and Palestine first formally recognised each other.
It was supposed to be the governing body of the state of Palestine, but the peace process never concluded and that state never effectively came into being.
Today, the PA is dominated by Fatah, the party of Mahmoud Abbas, Palestine ' s President.
In Gaza, Hamas has been in power ever since it defeated Fatah in the 2006 elections.
The secular and moderate Fatah is more acceptable to the West than Hamas, and the US and its allies are pushing for the former to come to power in Gaza too after the war.
Thus, Prime Minister Mohammad Shtayyeh and his Cabinet have resigned, and President Abbas is widely expected to pick as the next Prime Minister the US-educated Mohammad Mustafa.
Mustafa is an economist who worked at the World Bank, currently leads the Palestine Investment Fund, and is someone the West might find easier to work with.
As Shtayyeh said in his resignation statement,"I see that the next stage @ @ @ @ @ @ @ @ @ @ take into account the emerging reality in the Gaza Strip, the national unity talks, and the urgent need for an inter-Palestinian consensus based on a national basis, broad participation, unity of ranks, and the extension of the Palestinian Authority ' s sovereignty over the entire land of Palestine."He has been asked to stay on as \textbf{caretaker} PM till Abbas picks his successor.
Problems in this plan Advertisement Even if the PA is reformed, there are many roadblocks to it gaining power in Gaza.
First, of course, is the fact that the war is yet to end and Hamas is nowhere close to"total destruction", which Israeli Prime Minister Benjamin Netanyahu has said is his goal before he halts the offensive.
Second, Israel ' s plans for post-war Gaza do not envisage much autonomy for the region.
A future plan the Jewish nation has shared talks of Israel maintaining full control over security in a"demilitarised"Gaza, while civilian management and public order is in the @ @ @ @ @ @ @ @ @ @ would"not be identified with countries or entities that support terrorism and will not receive payment from them", as reported by BBC.
Third, Fatah is highly unpopular in Palestine, specially Gaza, and its people are unlikely to accept a government it heads.
Fatah is seen as corrupt, inefficient, and in cahoots with Israel in its brutal security crackdowns.
While Israel has accused Fatah of being too generous with the families of killed militants by giving them payouts, Palestinians feel Fatah has not done enough to advocate their cause.
Reasons this plan is being pushed Advertisement If the Gaza-under-PA plan seems doomed to fail, why is it being pushed?
Mainly because Israel ' s violence in Gaza is attracting louder criticism by the day, and this is one of the ways a ceasefire can possibly be achieved.
US President Joe Biden said on February 26 that he was hoping a ceasefire could be reached by next Monday.
In another signal that PA is being placated, there seemed to @ @ @ @ @ @ @ @ @ @ ( February 27 ).
US Treasury Secretary Janet Yellen said Israel had agreed to resume tax revenue transfers to the Palestinian Authority to fund basic services and bolster the West Bank economy, reported Reuters.
Fatah has consistently said it is unable to function effectively because Israel does not transfer to it the tax revenue it is supposed to.
It is being hoped that the reflow of money will prevent protests and riots in the West Bank.

% matched lemmas: caretaker
\end{textsample}
